%==============================================================================
% ABSTRACT
%==============================================================================
\chapter*{Abstract}
\addcontentsline{toc}{chapter}{Abstract}

\textbf{Keywords:} Retrieval-Augmented Generation, Model Context Protocol, SharePoint, Enterprise Knowledge Management, Large Language Models, Semantic Search, FAISS, Security, Hallucination.

\bigskip

Large organizations accumulate substantial volumes of technical and business documentation on enterprise platforms such as SharePoint, yet access to this knowledge remains fragmented and poorly suited to complex contextual queries. Large Language Models (LLMs) offer a natural language interface to enterprise knowledge, but their direct deployment is impeded by three structural limitations: the generation of unverified facts (hallucination), the absence of any mechanism for accessing private organizational documentation, and the impossibility of enforcing document-level access control.

This PFE report addresses the following research question: \textit{to what extent does a Retrieval-Augmented Generation (RAG) architecture grounded in the Model Context Protocol (MCP) enable more reliable, secure, and traceable access to enterprise knowledge stored in SharePoint, compared to a standalone LLM and a naive RAG implementation?}

We design, implement, and evaluate AntiGravity --- a modular RAG conversational assistant that leverages MCP as a standardized protocol layer between the AI reasoning engine and the SharePoint document corpus of the DNSI (Direction du Num\'{e}rique et des Syst\`{e}mes d'Information). The architecture separates a security-aware MCP Context Server, which manages a FAISS semantic index with ACL-filtered document access, from a conversational client that combines query classification, dense retrieval, cross-encoder re-ranking, and natural language generation via a locally deployed Mistral 7B Instruct model.

The evaluation follows a four-dimensional protocol: (1) retrieval quality (Precision@K, Recall@K, MRR, nDCG) on a domain-specific 50-item benchmark constructed from the DNSI corpus; (2) generation quality (faithfulness, groundedness, hallucination rate via RAGAS); (3) system performance (end-to-end latency, component-level profiling); and (4) user acceptance (SUS, TAM, task-based evaluation with DNSI employees).
% [NOTE FOR AUTHOR: Insert the actual quantitative results here once experiments are complete.
%  Replace this comment block with measured values:
%  - Hallucination rate: LLM alone (X%) vs. Naive RAG (Y%) vs. RAG+MCP (Z%)
%  - RAGAS Faithfulness: Cond. A / Cond. B / Cond. C
%  - MCP protocol overhead at P95 (ms)
%  - End-to-end latency P95 on CPU (s) and expected GPU estimate
%  - SUS score (mean, 95% CI), N participants
%  - TAM Perceived Usefulness vs. SharePoint baseline (mean, p-value)
%  - Task completion rate and mean time saving]

The primary scientific contribution is a formal comparative evaluation framework for RAG+MCP architectures applied to a private enterprise SharePoint corpus with ACL enforcement. The engineering contribution is a reference architecture for secure, permission-aware conversational AI in enterprise environments, deployable via Docker under corporate proxy constraints.

\bigskip

\noindent\rule{\textwidth}{0.4pt}

\bigskip

\textbf{R\'{e}sum\'{e} (French)}

\textbf{Mots-cl\'{e}s~:} G\'{e}n\'{e}ration Augment\'{e}e par R\'{e}cup\'{e}ration, Model Context Protocol, SharePoint, Gestion des Connaissances en Entreprise, Grands Mod\`{e}les de Langage, Recherche S\'{e}mantique, S\'{e}curit\'{e}, Hallucination.

\bigskip

Les grandes organisations accumulent de vastes quantit\'{e}s de documentation technique et m\'{e}tier sur des plateformes comme SharePoint, mais l'acc\`{e}s \`{a} ce patrimoine informationnel reste fragment\'{e} et insuffisant pour r\'{e}pondre \`{a} des requ\^{e}tes complexes et contextuelles. Bien que les grands mod\`{e}les de langage (LLMs) offrent une interface en langage naturel prometteuse, leur d\'{e}ploiement direct est entrav\'{e} par les ph\'{e}nom\`{e}nes d'hallucination, le manque de connaissance des donn\'{e}es priv\'{e}es, et l'absence de m\'{e}canismes natifs de contr\^{o}le d'acc\`{e}s et de tra\c{c}abilit\'{e} des sources.

Cette th\`{e}se con\c{c}oit, impl\'{e}mente et \'{e}value AntiGravity --- un assistant conversationnel RAG modulaire qui exploite le Model Context Protocol (MCP) comme couche de protocole standardis\'{e}e entre le moteur d'inf\'{e}rence IA et le corpus documentaire SharePoint de la DNSI. L'architecture s\'{e}pare un serveur de contexte MCP (index FAISS s\'{e}mantique avec filtrage ACL) d'un client conversationnel combinant classification de requ\^{e}te, recherche dense, re-classement par encodeur crois\'{e} et g\'{e}n\'{e}ration via Mistral 7B. L'\'{e}valuation suit un protocole en quatre dimensions (qualit\'{e} de r\'{e}cup\'{e}ration, qualit\'{e} de g\'{e}n\'{e}ration, performance syst\`{e}me, acceptation utilisateur)~; les r\'{e}sultats quantitatifs d\'{e}taill\'{e}s sont report\'{e}s au chapitre~\ref{chap:evaluation}.

% [NOTE FOR AUTHOR: Add key figures from the evaluation here once experiments are completed.]

\bigskip
\noindent\rule{\textwidth}{0.4pt}
